\section{Density Operator}
Prima di affrontare il discorso legato al \textbf{density operator} è importante approfondire quella che è una matrice di proiezione.

\begin{info}
\textbf{Projection Operator}\\
Una proiezione su uno spazio vettoriale \(\mathbb{V}\) è un operatore lineare \(P:\mathbb{V} \rightarrow \mathbb{V}\) tale per cui \(P^2=P\).
Se \(\mathbb{V}\) è uno spazio di Hilbert, può essere utilizzato il concetto di ortogonalità. Una proiezione \(P\) su uno spazio di Hilbert \(\mathbb{V}\) è detta ortogonale se soddisfa la seguente proprietà:
\[
<P x,y> = <x,P y>, \forall x,y \in \mathbb{V}
\]
Una proiezione sullo spazio di Hilbert che non è ortogonale è chiamata anche \textbf{proiezione obliqua}(citata solo 
per conoscenza e completezza).\\
\\
\textbf{Matrice di Proiezione}\\
La matrice di proiezione è una matrice quadrata che gode della proprietà elencata prima, ovvero: \(P^2=P\) (dove il prodotto tra due matrici è il prodotto riga per colonna).
Si definisce anche \textbf{matrice di proiezione ortogonale} se: \(P^2=P=P^T\) nel caso di spazi vettoriali reali, mentre nel caso di spazi vettoriali complessi si avrà \(P^2 = P = P^\dagger\).

Se la matrice di proiezione non è una matrice di proiezione ortogonale allora si dice \textbf{matrice di proiezione obliqua}.

Una conseguenza delle matrici di proiezione che sono matrici di proiezione ortogonale è che tutti gli autovalori associati saranno 1 e 0 con diversi gradi di molteplicità algebrica
\end{info}


L'\textbf{operatore di densità} è un modo alternativo di rappresentare gli stati e di formulare i postulati della meccanica quantistica. Tale strumento permette di poter esprimere meglio la gestione di alcuni scenari comuni nella meccanica quantistica. Una delle cose principali da capire è l'\textbf{ensembles of quantum states} (insiemi di stati quantistici), tale teoria ha come base quello presentato nella precedente sezione, ovvero, che uno stato quantistico non è detto sia puro, poichè non sarà possibile avere a propria disposizione lo stato di tutto il sistema, ma sarà misto, e quindi, confinato ad uno specifico sottositema. Dati quindi, gli stati \(\ket{\psi_i}\) in cui il sistema può trovarsi e \(p_i\) probabilità di trovarsi nell'\(i\)-esimo stato. Si chiamerà insieme di stati puri il duo: \(\{p_i,\ket{\psi_i}\}\). L'operatore di densità (o matrice di densità), sarà quindi definibile come:
\begin{gather*}
\rho = \sum_{i} p_i \ket{\psi_i} \bra{\psi_i}
\end{gather*}

Con tale definizione sarà possibile andare a riformulare anche quelli che sono i postulati stessi della meccanica quantistica, andando a spiegare anche quando tale notazione sia utile. 

\subsubsection{Legge di Shrodinger}
Quella che deriva dal quarto postulato della meccanica quantistica spiegando che l'evoluzione di un sistema quantistico è definito da matrici unitarie. Definendo quindi l'operatore unitario \(U\), considerando che inizialmente lo stato è \({p_i, \ket{\psi_i}}\), allora applicando l'evoluzione unitaria si avrà:
\(\ket{\psi_i} \rightarrow U \ket{\psi_i}\), dove la probabilità non è cambiata, quindi il nuovo stato sarà definito come: \(\{p_i, U \ket{\psi_i}\}\), quindi la probailità non cambia, andando ad analizzare quello che succede matematicamente, si avrà:

\begin{gather*}
\rho=\sum_{i} p_i \ket{\psi_i}\bra{\psi_i} \overset{U}{\longrightarrow} \sum_{i} p_i U \ket{\psi_i}\bra{\psi_i}U^\dagger = U \rho U^\dagger
\end{gather*}

\subsubsection{Postulato della misura}
Il postulato della misura è descrivibile mediente il linguaggio dell'operatore di densità. Supponiamo di effettuare una misurazione descritta da un operatore di misura \(M_m\). 

\begin{info}
\textbf{Proiettori degli autostati}\\
Dato un osservabile A ed i suoi autostati \(\{\ket{a_j}\}\) si introducono i proiettori su tali autostati come:
\begin{gather*}
M_j = \ket{a_j}\bra{a_j}
\end{gather*}

Un proiettore si dice \textbf{Hermitiano idempotente} quando soddisfa la proprietà tale per cui:
\begin{gather*}
P^2=P=P^\dagger
\end{gather*}

Queste sono esattamente le proprietà presentate precedentemente per l'operatore di densità ortogonale (Hermitiano + Idempotente).

Nei nostri casi, ad esempio, misurando l'osservabile \(Z\), i cui autostati sono \(\{\ket{0},\ket{1}\}\), si avrà che sarà possibile definire i seguenti due proiettori:
\begin{gather*}
M_0 = \ket{0}\bra{0} =
\left[\begin{smallmatrix}
1 & 0 \\
0 & 0
\end{smallmatrix}\right] \\
M_1 = \ket{1}\bra{1} =
\left[\begin{smallmatrix}
0 & 0 \\
0 & 1
\end{smallmatrix}\right]
\end{gather*}

Come visto anche in precedenza, essendo tale operatore ortonormale si avrà che i suoi autovalori saranno 0 ed 1.

Date queste nozioni, sarà quindi definibile come proiettore, un operatore unitario che aisce su un vettore di uno spazio vettoriale di Hilbert \(\mathcal{H}\) di dimensione \(N\), trasformandolo nella sua componente appartenente ad un sottospazio \(\hat{\mathcal{H}} \in \mathcal{H}\) di dimensione \(M\). Sia \(\{\ket{1}, \ket{2}, \dots, \ket{N}\}\) una base ortonormale di \(\mathcal{H}\) e sia \(\hat{\mathcal{H}}\) il sottospazio generato dai vettori \(\{\ket{1}, \ket{2}, \dots, \ket{M}\}\), il proiettore su tale stato è:
\begin{gather*}
P=\sum_{i=1}^{M} \ket{j}\bra{j}
\end{gather*}

Un operatore \(M\) è detto \textbf{normale} se: \(M M^\dagger = M^\dagger M\). Ma gli operatori unitari godono della proprietà per cui: \(U U^\dagger = I = U^\dagger U\). 

Gli operatori normali giocano un ruolo fondamentale per il teorema della decomposizione spettrale. Ogni operatore normale \(M\) su un vettore \(V\) è diagonale rispetto a qualche base ortogonale per \(V\). Di conseguenza, ogni operatore diagonalizzabile è normale. Questo significa che \(M\) può essere scritto come \(M=\sum_{i} \lambda_i \ket{i}\bra{i}\) dove i \(\lambda_i\) sono gli autovalori di \(M\), \(\{\ket{i}\}\) è una base ortogonale per \(V\), e ogni \(\ket{i}\) è un autovettore di \(M\) con autovalore \(\lambda_i\). Tali proiettori soddisfano la relazione di completezza, ovvero:
\begin{gather*}
\sum_{i} P_i = I
\end{gather*}
\end{info}

\begin{info}
L'uguaglianza segue dalla definizione di traccia,
$\operatorname{tr}(B)=\sum_k \bra{k}B\ket{k}$, calcolata in una base
ortonormale $\{\ket{k}\}$, e dalla relazione di completezza
$\sum_k \ket{k}\bra{k}=I$. Posto $A=M_m^\dagger M_m$, si ha
\begin{align*}
\operatorname{tr}\left(A\ket{\psi_i}\bra{\psi_i}\right)
  &= \sum_k \bra{k}A\ket{\psi_i}\braket{\psi_i | k} \\
  &= \sum_k \braket{\psi_i|k}\bra{k}A\ket{\psi_i} \\
  &= \bra{\psi_i}\Big(\sum_k \ket{k}\bra{k}\Big)A\ket{\psi_i}
   = \bra{\psi_i}A\ket{\psi_i},
\end{align*}
dove nel secondo passaggio si sono scambiati due scalari. Equivalentemente,
è un'applicazione della ciclicità della traccia,
$\operatorname{tr}(XY)=\operatorname{tr}(YX)$, con $X=A\ket{\psi_i}$ e
$Y=\bra{\psi_i}$, tenendo conto che la traccia di uno scalare è lo scalare
stesso.
\end{info}

Laprobabilità che, quindi, misurando l'osservabile \(A\) a partire dallo stato \(\ket{\psi}\) il sistema si porti nello stato \(\ket{a_j}\) è data da:
\begin{gather*}
p_j=Tr\left(M_j\rho\right)
\end{gather*}

La dimostrazione di tale proprietà è una conseguenza diretta:
\begin{gather*}
Tr\left(M_j\rho\right) = \sum_{k}\bra{a_k}M_j \rho \ket{a_k} = \sum_{k}\bra{a_k}\ket{a_j}\bra{a_j}\ket{\psi}\bra{\psi}\ket{a_k} =\\
=b_j b_j^* = |b_j|^2
\end{gather*}

Dove il termine \(\bra{a_k} \ket{a_j} \ne 0 \iff k=j\), e quindi, dato che siamo con una base ortonormale, \(\bra{a_k} \ket{a_k} = 1\), di conseguenza si avrà:
\begin{gather*}
\bra{a_k} \ket{\psi} \bra{\psi} \ket{a_k} = b_k b_k^* = |b_k|^2
\end{gather*}
Dato che stiamo andando a filtrare la \(k\)-esima componente del vettore \(\ket{\psi}\) mappata sul autovettore \(\ket{a_k}\). Ottenendo di conseguenza la relazione.

\subsubsection{Caratteristiche del density operator}
Il density operator gode di diverse proprietà:
\begin{itemize}
       \item \(\rho\) è Hermitiano, quindi: \(\rho = \rho^\dagger\)
       \item Traccia unitaria: \(Tr(\rho) = 1\)\\
       Tale proprietà è dimostrabile mediantei seguenti passaggi:
       \[
       Tr(\rho) = \sum_{k} \bra{k} \rho \ket{k} = \sum_{k} \sum_{i} p_i \bra{k}\ket{\psi_i}\bra{\psi_i}\ket{k} =
       \sum_{i}p_i \bra{\psi_i} \sum_{k}\ket{k}\bra{k}\ket{\psi_i} =
       \sum_{i}p_i = 1
       \]
       o equivalentemente:
       \[
       Tr(\rho) = Tr(\sum_{i}p_i\ket{\psi_i}\bra{\psi_i})=
       \sum_{i}p_i Tr(\ket{\psi_i}\bra{\psi_i})=
       \sum_{i} p_i = 1
       \]

       \item \(\rho\) è semidefinito positivo, ovvero:
       \[
       \bra{\psi}\rho\ket{\psi} \ge 0, \forall \ket{\psi}
       \]
       Tale proprietà discende dal fatto che gli autovalori di \(\rho\) sono tutti reali (essendo \(\rho\) Hermitiano) e non-negativi, quindi dalla condizione: \(Tr(\rho) = 1, 0 \le \lambda_i \le 1\).

       La dimostrazione di tale proprietà discende dalla definizione data precedentemente:
       \[
       \bra{\varphi}\rho\ket{\varphi} = \sum_{i}p_i\braket{\varphi | \psi_i}\braket{\psi_i | \varphi}=\sum_{i} p_i |\braket{\varphi | \psi_i}|^2 \ge 0
       \]
       Dato che \(p_i\) essendo una probabilità non potrà essere mai zero. Lo stesso valer per il modulo del numero complesso ottenuto.

       Se supponiamo che \(\rho\) è un operatore che soddisfa le condizioni di traccia e di positività, si avrà che:
       \[
       \rho=\sum_{j}\lambda_j \ket{j}\bra{j}
       \]
       dove i vettori \(\ket{j}\) sono ortogonali e i \(\lambda_i\) sono autovalori reali e non-negativi di \(\rho\).
\end{itemize}

Per uno stato puro \(\ket{x}\), l'operatore di densità \(\rho_x^X=\ket{x}\bra{x}\) ha una forma perticolarmente semplice. Se consideriamo che \(\ket{x}\) è l'i-esimo elemento di una base ortonormale, allora \(\ket{x}\) è il vettore della base con tutti 0 ed un solo 1 nella i-esima posizione. Di conseguenza, il prodotto \(\ket{x}\bra{x}\) sarà uguale ad una matrice composta da tutti 0 e con un solo 1 presente sulla diagonale alla posizione \(ii\). Segue quindi la proprietà che l'operatore di densità di uno stato puro è un proiettore di quest'ultimo, ovvero: \(\rho_x^X \rho_x^X = \rho_x^X\). Contrariamente, ogni operatore di densità che è un proiettore corrisponde ad uno stato puro: operatori di proiezione hanno solo 0 ed 1 come autovalori, e per ottenere la traccia uguale ad 1, l'operatore di densità avrà un solo ed unico autovettore, che corrisponderà allo stato puro di quell'operatore.

Nel caso in cui mi trovi in degli stati misti, invece, l'operatore di densità, sulla diagonale, non avrà tutti 1 e 0, questo si dice che lo stato si un misto di diversi stati puri nell'insieme per \(\rho\). Uno stato puro, quindi, soddisfa il criterio per cui \(tr(\rho^2)=1\) questo perchè avendo un solo 1 sulla diagonale, il risultato è "ovvio". Mentre nel caso in cui ci sia uno stato misto, si avrà che \(tr(\rho^2) < 1\), questo perchè la somma degli autovalori \(\lambda_j\) sulla diagonale è uguale ad 1, ma questo impone che tali valori siano \(0 \le \lambda_j \le 1\), di conseguenza si avrebbe che la traccia sia definita come: \(Tr(\rho^2)=\sum_{i}\lambda_i^2 < 1\). Quindi, dato un operatore di densità \(\rho\), la decomposizione \(\rho=\sum_{i} p_i \ket{\psi_i}\bra{\psi_i}\) non è unica. 


% % Roba Vecchia
% si definisce come:
% \[
% \rho = \sum_{i} p_i \ket{\psi_i}\bra{\psi_i}
% \]

% Dove il termine \(\rho\) incapsula all'interno della sua struttura sia la probabilità dello stato puro \(p_i\) che lo stato puro stesso \(\ket{\psi_i}\).
% Per un sistema, invece, dove non andiamo a considerare una sua restrizione, ma un suo stato puro, si avrà:
% \[
% \rho = \ket{\psi}\bra{\psi}
% \]
% E tutti i postulati della meccanica quantistica possono essere espressi in termini di operatori di densità


